% !TeX root = ../../tesis.tex

% =============================================================================
%  B.8 — INFRAESTRUCTURA: DOCKER Y AUTOMATIZACIÓN
% =============================================================================
\section{Infraestructura: Contenedores y Automatización}
\label{anx:vft-infraestructura}

El despliegue de VFTModel se basa en tres archivos de configuración que permiten
levantar el servicio tanto en entorno local como en producción sin modificar el
código fuente.

\subsection*{Dependencias Python}

El archivo \texttt{requirements.txt} agrupa las dependencias en tres categorías
funcionales:

\begin{lstlisting}[style=estiloPython,
    caption={Dependencias de VFTModel (\texttt{requirements.txt}).},
    label=lst:vft-requirements]
# Web y asincronia
fastapi==0.110.0    uvicorn==0.29.0
httpx==0.27.0       pydantic==2.6.4
python-dotenv==1.0.1

# Analisis espacial y grafos (core)
geopandas==0.14.3   networkx==3.2.1
momepy==0.7.0       shapely==2.0.3
numpy>=1.24.0       scipy>=1.10.0

# Visualizacion
matplotlib==3.8.3
\end{lstlisting}

El uso de \textsc{gdal} a nivel de sistema operativo es un requisito previo
ineludible para GeoPandas y Fiona:

\begin{lstlisting}[language=bash,
    caption={Instalación de GDAL a nivel de sistema operativo.},
    label=lst:vft-gdal]
# Linux (Debian/Ubuntu)
sudo apt install gdal-bin libgdal-dev
# macOS
brew install gdal
\end{lstlisting}

\subsection*{Dockerfile multietapa}

La imagen Docker utiliza una construcción multietapa (
\termino{multi-stage build}) para minimizar el tamaño de la imagen de
producción: el \textit{stage} \texttt{builder} compila e instala las
dependencias con todas las herramientas de desarrollo, y el \textit{stage}
\texttt{runtime} solo copia los paquetes ya instalados y el código fuente.

\begin{lstlisting}[language=bash,
    caption={Dockerfile multietapa de VFTModel.},
    label=lst:vft-dockerfile]
# Stage 1: Builder — instala dependencias con compilador
FROM python:3.12-slim-bookworm AS builder
RUN apt-get update && apt-get install -y --no-install-recommends \
    libgdal-dev gdal-bin libgeos-dev build-essential \
    && rm -rf /var/lib/apt/lists/*
WORKDIR /app
COPY requirements.txt .
RUN pip install --no-cache-dir -r requirements.txt

# Stage 2: Runtime — imagen minima sin compilador
FROM python:3.12-slim-bookworm AS runtime
RUN apt-get update && apt-get install -y --no-install-recommends \
    libgdal32 libgeos-c1v5 && rm -rf /var/lib/apt/lists/*
WORKDIR /app
COPY --from=builder /usr/local/lib/python3.12/site-packages \
                    /usr/local/lib/python3.12/site-packages
COPY src/ ./src/
ENV APIMETRO_URL=https://apimetro.dev/movilidad
EXPOSE 8000
CMD ["uvicorn", "src.api.main:app", "--host", "0.0.0.0", "--port", "8000"]
\end{lstlisting}

La imagen base \texttt{python:3.12-slim-bookworm} está fijada a la distribución
\textit{Bookworm} (Debian 12): la distribución siguiente (\textit{Trixie}) no
incluye el paquete \texttt{libgdal32} en sus repositorios estables al momento de
escritura del presente trabajo.

\subsection*{Makefile: automatización de tareas frecuentes}

\begin{lstlisting}[language=bash,
    caption={Comandos del \texttt{Makefile} de VFTModel.},
    label=lst:vft-makefile]
make install           # pip install -r requirements.txt
make run               # uvicorn LOCAL  (puerto 8000, .env.local)
make run-dev           # uvicorn DEV    (apimetro.dev, .env.dev)
make docker-build      # construye la imagen Docker
make docker-run        # Docker en modo DEV (apimetro.dev)
make docker-run-local  # Docker LOCAL — agrega --add-host para CDMX local
make test              # pytest (requiere servidor activo + apimetro)
make export-geojson    # exporta capas GeoJSON para Tableau/QGIS
\end{lstlisting}

El \termino{override} de puerto se aplica con \texttt{PORT=<n>}:
\texttt{make run PORT=9000}.
